\section{Flows,Cuts,Connectivity}
\begin{definition}
    A \textbf{directed graph} is a pair $(V,E)$ with $E\subseteq (V\times V)\backslash\{(x,x):x\in V\}$. Write $\delta_G^+(v)$ for edges starting in $v$ and $\delta_G ^-(v)$ for edges ending in $v$.
\end{definition}
\vspace{0.5cm}
\begin{definition}
    A \textbf{flow network} is a tuple $(G,s,t,c)$ where $G$ is a graph, $s,t\in V(G),s\neq t$ and $c:E\to\mathbb{R}_0^+$.
    \begin{itemize}
        \item s is called a \textbf{source}
        \item t is called a \textbf{target}
        \item $\forall e\in E(G):c(e)$ is the \textbf{capacity of e}
    \end{itemize}
    A \textbf{flow} in a flow network is a function $f:E\to\mathbb{R}^+_0$ s.t.
    \begin{itemize}
        \item $f(e)\leq c(e)$ for all $e\in E(G)$
        \item $\forall v\notin\{s,t\},\sum\limits_{e\in\delta^-(v)} f(e)=\sum\limits_{e\in\delta^+(v)}f(e)$ (i.e. what flows in flows out)
        \item $\forall e\in\delta^-(s)\cup\delta^+(t):f(e)=0$ (i.e. nothing flows into a source or out of a target)
    \end{itemize}
    The \textbf{value} of a flow $f$ is $\operatorname{val}(f)=\sum\limits_{e\in\delta^-(v)} f(e)=\sum\limits_{e\in\delta^+(v)}f(e)$
\end{definition}
\begin{remark}
    For $v\in V$, write $f(\delta^+(v))=\sum\limits_{e\in\delta^+(v)} f(e) \text{ ; }f(\delta^-(v))=\sum\limits_{e\in\delta^+(v)}f(e)$
\end{remark}
\vspace{0.5cm}
\begin{definition}
    A \textbf{cut} in a graph/directed graph $G$ is a partition of $V(G)$ into two non-empty parts $A$ and $B$ (i.e. $V(G)=A\stackrel{.}{\cup} B$ and $A,B\neq\emptyset$). For $s,t\in V(G),s\neq t \implies (A,B)$ is a $(s,t)$-cut if $s\in A,t\in B$.\\
    If $(S,T)$ is a $(s,t)$-cut, write $\delta^+(S)=\{e\in E(G):e$ start in $S$ ends in $T\}=\delta^-(T)$ and $\delta ^-(S)=\{e\in E(G): e$ starts in $T$ and ends in $S\}=\delta^+(T)$.\\
    The \textbf{capacity} of the cut is $c(S,T)=\sum\limits_{e\in \delta^+(S)}c(e)=c(\delta^+(S))=c(\delta^-(T))$.
\end{definition}
\vspace{0.5cm}
\begin{theorem}[max flow-min cut theorem]
    Let $(G,s,t,c)$ be a flow network. Then $\max\{\operatorname{val}(f):f\text{ flow}\}=\min\{c(S,T):(S,T) \text{ is a } (s,t)\text{-cut}\}$.
\end{theorem}
\vspace{0.5cm}
\begin{lemma}
    Let $f$ be a flow, $(S,T)$ be a $(s,t)$-cut. Then $\operatorname{val}(f)=f(\delta^+(S))-f(\delta^-(S))$.
\end{lemma}
\begin{proof}
    \hfill \\
    Claim: $\forall A\subset V:f(\delta^+(A))-f(\delta^-(A))=\sum\limits_{v\in A} f(\delta^+(v))-f(\delta^-(v))=(\ast)$\\
    The sum $f(\delta^+(A))-f(\delta^-(A))$ does not account for edges $e\in E(G)$ with start and end in $A$. On the right side, $f(e)$ is counted twice, once with positiv and once with negativ sign, so the sum vanishes.\\
    Then $\operatorname{val}(f)=f(\delta^-(t))\stackrel{(1)}{=}\sum\limits_{v\in T}f(\delta^-(v))-f(\delta^+(v))\stackrel{(\ast)}{=}f(\delta^-(T))-f(\delta^+(T))=f(\delta^+(S))-f(\delta^-(S))$\\
    $(1)$ holds because $f(\delta^+(t))=0$ ($t$ is a sink) and $\forall v\notin \{s,t\}:f(\delta^-(v))-f(\delta^+(v))=0$ (by definition).
\end{proof}
\begin{corollary}
    Let $f$ be a flow, $(S,T)$ an $(s,t)$-cut.\\
    $\val(f)\leq c(\delta^+(S))$\\
    Thus $\max\{\val(f): f \text{ flow}\}\leq \min\{c(S,T):(S,T) \  (s,t)\text{-cut}\}$.
\end{corollary}
\begin{proof}
    \hfill
    $\max$ flow $\leq\min$ cut due to Corollary $(4.7)$\\
    Let $f$ maximum flow (i.e. a flow for which $\val(f)$ is maximal). Let $c':E(G)\to \mathbb{R}_0^+$ where $c'(e)=c(e)-f(e)$. Let $S\subset V(G)$ be defined as the set of reachable vertices according to the following rules:\\
    \begin{itemize}
        \item $s\in S$
        \item $x\in V(G)$ is reachable if there is a walk (in the underlying directed graph) from $s$ to $x$ traversing $e_1,...,e_k$ s.t. 
    \end{itemize}
    $c(e_i)>0$ ($f(e_i)<c(e_i)$) if $e_i$ is traversed in the correct direction and $c'(e_i)<c(e_i)$ ($f(e_i)>0$) if $e_i$ is traversed in the opposit direction.\\
    If $S=V(G)$, we can improve $f$ as follows:\\
    Let $e_1,...e_k$ be a walk from $s$ that makes $t$ reachable, and $e_i$ is traversed in orientation $\sigma_i\in\{\pm 1\}$, then $\tilde{f}(e_i)=f(e_i)+\sigma_i\eps$, for $i\in\{1,...,k\}$ and $\tilde{f}(e_i)=f(e_i)$ otherwise. This defines a flow $\tilde{f}$ with $\val(\tilde{f})>\val(f)$ for $\eps>0$ sufficently small $\lightning$\\
    Thus $t\notin S$. Set $T=V\setminus S\neq\emptyset$
    \begin{itemize}
        \item $\forall e\in\sigma^+(S)\quad c'(e)=0$, i.e. $f(e)=c(e)$\\
        otherwise we would get a path from $s$ to the start $x$ of $e$, which makes $x$ reachable, and extends by $e\lightning$
        \item $\forall e\in \delta^-(S)\quad c'(e)=c(e)\iff f(e)=0$ by the same argument
        \item $c(S,T)=\sum\limits_{e\in\delta^+(S)}c(e)=\sum\limits_{e\in\delta^+(S)}f(e)-\sum\limits_{e\in\delta^-(S)}f(e)=f(\delta^+(S))-f(\delta^-(S))\stackrel{(4.6)}{=}\val(f)\implies\val(f)\geq c(S,T)\geq \min(s,t)$-cut
    \end{itemize}
\end{proof}

\subsection{Connectivity}
